\subsection{Αναλύτικη περιγραφή ενός πειράματος και πως φαίνεται στην ιστοσελίδα}


Στην εικόνα \ref{fig:defaultSite},φαίνεται η αρχική κατάσταση της ιστοσελίδας που αλλελπιδρά ο χρήστης.
Είναι φτιαγμένη με σκοπό να μην χρειάζονται ειδικές γνώσεις για να χειριστεί κάποιος την εφαρμογή.

Στο κάτω μέρος υπάρχει το κομμάτι που ο χρήστης δηλώνει υγρασία και θερμοκρασία,αλλά όπως αναφέρθηκε και προηγούμενης,δεν χρησιμοποιήθηκε στα προκειμένα πειράματα.
Δίνεται πρόσβαση στο url που κατευθείνει τον χρήστη στην grafana όπου υπάρχει η οπτικοποίηση σε ζωντανό χρόνο.
Επίσης,υπάρχει κομπί για μία περιγραφή της διαδικασίας του πειράματος και δυνατότητα να αποσταλούν τα σημερίνα log του server στον χρήστη.

\begin{figure}[!htbp]
    \centering
    \includegraphics[width=0.8\textwidth]{}
    \caption{Αρχική μορφή της ιστοσελίδας}
    \label{fig:defaultSite}
\end{figure}

Αφού ο χρήστης έχει κλείσει τα παράθυρα και την πόρτα και έχει περιμένει κάποια λεπτά, ώστε η τιμή VOC να αυξηθεί από την κατάσταση που το δωμάτια έχει ανοιχτά το παράθυρο και την μπαλκονόπορτα,
που θα αρχίζουμε να τα λέμε και τα δύο παράθυρα,γιατί και τα δύο επιτελούν τον ρόλο του να αερίζουν το δωμάτιο με εξωτερικό αέρα.

Αφού παρατηρηθεί από την grafana ότι οι τιμές VOC έχουν σταθεροποιηθεί σχετικά,δηλώνει ότι θέλει να ξεκινήσει καινούργιο πείραμα,πατώντας το κουμπί \textbf{Ξεκίνα καινούργιο πείραμα}.
Ο server αποθηκεύει στο collection  \texttt{UserInput} της βάσης δεδομένων,document με \texttt{userInputCategory="ExperimentState"} και \texttt{uexperimentState="StartingExperiment"} και την χρονική στιγμή της αποθήκευσης του document.
Έπειτα εμφανίζεται η παρακάτω μορφή της ιστοσέλίδας \ref{fig:afterInsertingStart}. 


\begin{figure}[!htbp]
    \centering
    \includegraphics[width=0.8\textwidth]{}
    \caption{Η ιστοσελίδα αφού έχει ξεκίνησει πείραμα}
    \label{fig:afterInsertingStart}
\end{figure}

Αφού περιμένουμε τουλάχιστον 1 με 2 λεπτά,για να υπάρχει περιθώριο δεδομένων ο χρήστης πατάει το κουμπί \textbf{},
για να προχωρήσει στην σελίδα που φαίνεται στην εικόνα \ref{fig:InsertingSourcePollutantDetails}.

\begin{figure}[!htbp]
    \centering
    \includegraphics[width=0.8\textwidth]{}
    \caption{Η σελίδα που δηλώνονται οι λεπτομέρειες του πειράματος}
    \label{fig:InsertingSourcePollutantDetails}
\end{figure}

\FloatBarrier 


\subsubsection{Λεπτομέρειες τις σελίδας εισαγωγής πηγής ρύπων }




\paragraph{Τύπος Ρύπου,Αντικείμενο προέλευσης,ποσότητα που χρησιμοποείται}

\paragraph{Δωμάτιο}

\paragraph{Θέση πηγής ρύπου}

\paragraph{Συνθήκες του δωματίο που ισχύουν}

\begin{itemize}
  \item
  \item 
  \item 
\end{itemize}
\subsubsection{Εισαγωγή πηγής ρύπου}

\begin{figure}[!htbp]
    \centering
    \includegraphics[width=0.8\textwidth]{}
    \caption{Η σελίδα αφότου δηλώθηκε η είσοδος πηγής ρήπου}
    \label{fig:sourcePollutantInserted}
\end{figure}


\FloatBarrier 
